\begin{textsample}{POS Dim 2 – rs – Score 33.00 – rs/rs\_em\_86.txt}  \label{ex:f2_pos_123}
6.4 FORMAÇÃO PARA O TRABALHO : RECONHECIMENTO DE SABERES E PRÁTICAS A Pesquisa Nacional por Amostra de Domicílios ( PNAD ), revela que o Brasil tem 47,3 \textbf{milhões} de jovens, de 15 a 29 anos de idade.
Desse total, 13,5 \% estavam ocupados e estudando ; 28,6 \% não estavam ocupados, porém estudavam ; 34,9 \% estavam ocupados e não estudavam.
Finalmente, 23 \% não estavam ocupados nem estudando.
Assim, quase metade deles, 48,4 \%, ou seja, 22,8 \textbf{milhões} de jovens brasileiros estão “ ocupados ” economicamente.
Já na faixa dos 14 aos 17 anos, esse índice alcança 46,36 \% da população jovem.
Nessa jornada, no ambiente das ocupações econômicas, os jovens acumulam saberes e práticas do mundo do trabalho, conforme já descreve a Resolução CNE/CP nº. 01/2021, em seu \textbf{artigo} 14 : “ (... ) o desenvolvimento de saberes instrumentais relacionados ao mundo do trabalho, na perspectiva da geração de trabalho e renda ” ( BRASIL, 2021e ).
Em atenção a esse cenário, entende -se indispensável reconhecer a formação no trabalho também no percurso do jovem que \textbf{cursa} algum Itinerário FTP.
É necessário incluir a possibilidade do reconhecimento de saberes profissionais, advindos da formação no trabalho, ao estudante matriculado em Itinerário FTP, o que é de responsabilidade da \textbf{Instituição} Educacional que oferece tal modalidade de Educação Profissional e Tecnológica.
Isso constitui tão importante aspecto da Educação Profissional e Tecnológica, que a Diretriz Nacional da EPT, através da Resolução CNE/CP nº. 01/2021 lhe \textbf{destinou} um \textbf{capítulo}, denominado “ \textbf{Capítulo} XV do reconhecimento de saberes e competências ”, que estabelece : Art. 47.
Os saberes adquiridos na Educação Profissional e Tecnológica e no trabalho podem ser reconhecidos mediante processo formal de avaliação e reconhecimento de saberes e competências profissionais - \textbf{Certificação} Profissional para fins de exercício profissional e de \textbf{prosseguimento} ou conclusão de estudos, em consonância com o \textbf{art}. 41 da Lei nº 9.394/1996. (... ) § 1º A \textbf{certificação} profissional abrange a avaliação do \textbf{itinerário} profissional e social do estudante, que inclui estudos não formais e experiência no trabalho ( saber informal ), bem como a orientação para continuidade de estudos, segundo \textbf{itinerários} formativos coerentes com os históricos profissionais dos cidadãos, para valorização da experiência extraescolar. ( BRASIL, 2021e ).
Por oportuno, sublinha -se a possibilidade de incluir o aproveitamento de atividades como a realização de estágios, conforme a Lei nº. 11.788/2008, que dispõe sobre o estágio de estudantes, bem como aproveitar a realização de \textbf{cursos} de formação técnico-profissional em \textbf{entidades} qualificadas.
Os contratos de Aprendizagem, de acordo com o Decreto nº. 9.579/2018, que consolidam normas da Aprendizagem, também são válidos nesse sentido e estão \textbf{previstos} na Resolução CNE/CP nº. 01/2021, no \textbf{Artigo} 5º : § 6º Os \textbf{itinerários} formativos profissionais devem possibilitar um contínuo e articulado aproveitamento de estudos e de experiências profissionais devidamente avaliadas, reconhecidas e certificadas por \textbf{instituições} e redes de Educação Profissional e Tecnológica, criadas nos termos da \textbf{legislação} \textbf{vigente}. § 7º Os \textbf{itinerários} formativos profissionais podem ocorrer dentro de um \textbf{curso}, de uma área tecnológica ou de um eixo tecnológico, de modo a favorecer a verticalização da formação na Educação Profissional e Tecnológica, possibilitando, quando possível, diferentes percursos formativos, incluindo programas de aprendizagem profissional, observada a \textbf{legislação} trabalhista pertinente. ( BRASIL, 2021e ).
Na mesma lógica do olhar sobre a ocupação econômica dos jovens, existe ainda um diversificado conjunto de estudos adicionais que se fazem presentes em suas rotinas.
Pode -se exemplificar com a realização de \textbf{curso} de línguas estrangeiras, de intercâmbios, de \textbf{cursinhos} e \textbf{cursos} de extensão em diversos campos de conhecimento e de interesse dos estudantes.
O cenário da diversificação de \textbf{oferta} de formações possíveis, especialmente, por meios eletrônicos, se acentuou no cenário pandêmico desde 2019, e a facilidade com que as juventudes “ navegam ” nas novas tecnologias, também amplia mais ainda o portfólio dessas possibilidades.
Esses saberes podem e devem ser considerados nos \textbf{Itinerários} Formativos do RCGEM.
Os saberes dos estudantes, trazidos além do \textbf{currículo} formal, são tema do"\textbf{Capítulo} XIV - do Aproveitamento de Estudos ”, da Diretriz da EPT supracitada, e \textbf{prevê} que os estudos, conhecimentos e experiências anteriores, inclusive no trabalho, devam ser reconhecidos, com vistas ao \textbf{prosseguimento}, desde que diretamente relacionados com o perfil profissional de conclusão da respectiva \textbf{qualificação} profissional ou \textbf{habilitação} profissional \textbf{técnica} ou tecnológica.
Novamente, a \textbf{Instituição} que \textbf{ofertar} algum Itinerário de FTP, elabora e submete ao CEEd-RS o \textbf{detalhamento} de como se dará tal aproveitamento, conforme o \textbf{previsto} no \textbf{artigo} 46 : I – em \textbf{qualificações} profissionais técnicas e unidades curriculares, etapas ou módulos de \textbf{cursos} \textbf{técnicos} ou de Educação Profissional e Tecnológica de \textbf{Graduação} regularmente concluídos em outros \textbf{cursos} ; II - em \textbf{cursos} \textbf{destinados} à \textbf{qualificação} profissional, incluída a formação inicial, mediante avaliação, reconhecimento e \textbf{certificação} do estudante, para fins de \textbf{prosseguimento} ou conclusão de estudos ; III - em outros \textbf{cursos} e programas de Educação Profissional e Tecnológica, inclusive no trabalho, por outros meios formais, não formais ou informais, ou até mesmo em outros \textbf{cursos} superiores de \textbf{graduação}, sempre mediante avaliação do estudante ; IV - por reconhecimento, em processos formais de \textbf{certificação} profissional, realizado em \textbf{instituição} devidamente credenciada pelo órgão normativo do respectivo sistema de ensino ou no âmbito de sistemas nacionais de \textbf{certificação} profissional de pessoas. ( BRASIL, 2021e ).
No que tange à \textbf{certificação} profissional, a \textbf{Portaria} MEC nº. 24/2021 dispõe sobre o Sistema Nacional de Reconhecimento e \textbf{Certificação} de Saberes e Competências Profissionais ( Re-Saber ).
Embora as possibilidades e orientações do Re-Saber se apliquem aos maiores de 18 anos, pode ser também uma alternativa aplicável ao Itinerário Formativo FTP, conforme aponta seu \textbf{art}. 2º : Art. 2º O processo de \textbf{certificação} profissional, no âmbito do Re-Saber, constitui -se em um conjunto articulado de ações de natureza educativa para : I - a sistematização de saberes e competências que possibilite a elaboração do processo de \textbf{certificação} profissional ; II - o desenvolvimento de metodologias que permitam identificar, avaliar e reconhecer saberes e competências que habilitem para o exercício profissional ou para a conclusão ou \textbf{prosseguimento} de estudos ; III - o atendimento às demandas de \textbf{certificação} profissional correspondentes aos \textbf{cursos} de \textbf{qualificação} profissional, \textbf{técnicos} de nível médio, de \textbf{especialização} \textbf{técnica} e superiores de tecnologia ; (... ) V - o estímulo à inclusão socioprodutiva e ao aumento das possibilidades de inserção profissional dos \textbf{trabalhadores} certificados ; VI - o incentivo à continuidade de estudos para a \textbf{elevação} da \textbf{escolaridade}, sempre que possível ; (... ). ( BRASIL, 2021c ).
Ainda com vistas a subsidiar as \textbf{Instituições} Escolares na elaboração de suas \textbf{diretrizes} e estratégias para aproveitamento de saberes e experiências no Itinerário FTP, é pertinente a \textbf{observância} da Resolução MEC/CNE/CEB nº. 3/2018, que \textbf{atualiza} as DCNEM : Art. 18.
Para efeito de cumprimento das exigências curriculares do ensino médio, os sistemas de ensino devem estabelecer critérios para reconhecer as competências dos estudantes, tanto da formação geral básica quanto dos \textbf{itinerários} formativos do \textbf{currículo} (... ) (... ) Parágrafo único.
No âmbito do \textbf{itinerário} de formação \textbf{técnica} e profissional, as \textbf{instituições} e redes de ensino devem realizar processo de avaliação, reconhecimento e \textbf{certificação} de saberes e competências adquiridos na educação profissional, inclusive no trabalho, para fins de \textbf{prosseguimento} ou conclusão de estudos nos termos do \textbf{art}. 41 da LDB, conferindo aos aprovados um diploma, no caso de \textbf{habilitação} \textbf{técnica} de nível médio, ou certificado idêntico ao de \textbf{curso} correspondente, no caso de curso(s) de \textbf{qualificação} profissional. ( BRASIL, 2018j ).

% matched lemmas: art, artigo, atualizar, capítulo, certificação, currículo, cursar, curso, destinar, detalhamento, diretriz, elevação, entidade, escolaridade, especialização, graduação, habilitação, instituição, itinerário, legislação, milhão, observância, oferta, ofertar, portaria, prever, prosseguimento, qualificação, trabalhador, técnico, vigente
\end{textsample}
